\documentclass[11pt,a4paper]{article}
\usepackage[margin=27mm]{geometry}
\usepackage{amsmath,amssymb,amsthm,fontspec,unicode-math,booktabs,longtable,hyperref,fvextra}
\setmainfont{STIXGeneral}\setmathfont{STIX Two Math}\setmonofont{Menlo}[Scale=MatchLowercase]
\setlength{\emergencystretch}{3em}
\newtheorem{theorem}{Theorem}\newtheorem{lemma}[theorem]{Lemma}\newtheorem{corollary}[theorem]{Corollary}\newtheorem{proposition}[theorem]{Proposition}
\newcommand{\Aut}{\operatorname{Aut}}\newcommand{\Orb}{\operatorname{Orb}}\newcommand{\supp}{\operatorname{supp}}
\hypersetup{colorlinks=true,linkcolor=blue,urlcolor=blue}

\title{Diagonal twistings and finite tree diagrams}
\author{Human Compression Project}\date{}
\begin{document}\maketitle
\begin{abstract}
We derive the two diagonal constructions and the finite tree-diagram theorem of Piedade--Tranchida from the common theory of orbit intersections. We give all maximal parabolics, exact diagram labels, invariant-quotient variants and the identifications with classical polytope extensions. The input is the complete twelve-page author edition arXiv:2608.23223v1. The general existence question for arbitrary diagrams is not answered here.
\end{abstract}
The common paper [4] supplies the coset and element-orbit results
(Lemma~2 and Theorem~3), the faithful regular construction
(Corollary~22), the quotient reduction (Lemma~23), the orbital
presentation (Proposition~24) and the exact doubled label (Lemma~25).
All geometric uses below have finite type sets. The quotient conclusions
begin with faithful regular inputs and retain the separate hypothesis
that the quotient maximal-parabolic geometry is regular; its action
kernel is removed as described below.

\section{The diagonal data}
Let $\beta=(B,(B_i)_{i\in I})$ be a finite-rank coset incidence system and fix $j\in I$. Assume that $X=B/B_j$ is finite and $|X|>1$. A $j$-diagonal is an unordered pair of distinct elements of $X$. Assign a label in $\{2,3,\ldots\}\cup\{\infty\}$ to each $B$-orbit of these pairs. Write $D$ for the resulting Coxeter graph on $X$, omitting edges with label $2$, and $A=W(D)$ for its Coxeter group. The notation $A(Z)$ means the subgroup generated by the Coxeter generators indexed by $Z\subseteq X$.

Every pair orbit has a representative $\{B_j,\tau B_j\}$ with $\tau\notin B_j$. The common orbital-presentation proposition proves that two such representatives agree exactly when
\[
\tau'\in B_j\tau B_j\ \cup\ B_j\tau^{-1}B_j.
\]
Thus the labels depend on inverse-paired double cosets, not just on powers of a convenient element. Left multiplication preserves these labels and acts on $D$, and hence on $W(D)$ by $a_F\mapsto a_{bF}$. The kernel of the action on $D$ is $\bigcap_{b\in B}bB_jb^{-1}$: an element fixes all left cosets exactly when each conjugate belongs to $B_j$. Standard Coxeter generators are distinct, so the induced action on $W(D)$ has the same kernel. Neither construction below requires this action to be faithful.

The action permutes the standard maximal parabolics transitively. If $B$ is flag-transitive on $\beta$, the common element-orbit theorem gives
\begin{equation}\label{eq:admissibility}
B_HB_j/B_j\ \cap\ B_KB_j/B_j=B_{H\cup K}B_j/B_j.
\end{equation}
Equivalently, with $O(M)=B_{I\setminus M}B_j/B_j$, one has $O(M)\cap O(N)=O(M\cap N)$. Without flag-transitivity this identity is precisely the remaining admissibility requirement; it is not automatic for an arbitrary incidence system.

If the regular input is supplied by a flag-transitive but unfaithful coset representation $B$, let $K=B_I$. Thinness and chamber connectivity show that $K$ is the normal action kernel. It lies in every $B_i$ and acts trivially on $X=B/B_j$. Thus the action on the Coxeter group also factors through $B/K$. Passing to $B/K$ preserves every parabolic intersection, every diagonal orbit and its label. The kernel $1\rtimes K$ is normal in $A\rtimes B$ and lies in all the twisting parabolics. The resulting coset geometry is therefore the same as the one for $A\rtimes(B/K)$. This reduces a flag-transitive unfaithful model to the faithful model below; the full-group assertions are made only for the faithful model. No replacement of an arbitrary non-flag-transitive group action by a regular one is being assumed.

\section{One orbit, one new type}
Assume now that $\beta$ is a regular hypertope in its faithful model. Let $x_0=B_j$ and let $b_i$ be its distinguished involutions. The classical Coxeter parabolic theorem makes the standard geometry of $A$ a regular hypertope. The common regular semidirect theorem applies by (\ref{eq:admissibility}) and gives the following complete construction.

\begin{theorem}\label{thm:one}
The group $G=W(D)\rtimes B$ is the full type-preserving automorphism group of a regular hypertope with types $\{*\}\sqcup I$, distinguished generators $a=a_{x_0}$ and $b_i$, and maximal parabolics
\begin{equation}\label{eq:onepars}
G_*=B,\qquad G_i=A(B_ix_0)\rtimes B_i\quad(i\in I).
\end{equation}
Its Coxeter diagram is the diagram of $\beta$ with a new node $*$ joined only to $j$, by the label $2m$, where
\[
m=\lambda([x_0,b_jx_0]).
\]
The label is exact, including when $m=\infty$. The geometry is finite exactly when both $B$ and $W(D)$ are finite. If $M\leq W(D)$ is invariant under $B$, and the standard quotient coset geometry for $W(D)/\langle\!\langle M\rangle\!\rangle$ is a regular hypertope, its corresponding twisting is also a regular hypertope.
\end{theorem}
\begin{proof}
There is just one orbit on the Coxeter types. The displayed maximal parabolics are exactly the common theorem's formulas; in particular the new parabolic is $A(\varnothing)\rtimes B$, not $A_{x_0}\rtimes B$. Their flag stabilizers, thinness, residual connectivity and full automorphism group are supplied by that theorem. Transitivity on $X$ also shows directly that $a,b_i$ generate $G$.

For $i\ne j$, the generator $b_i$ belongs to $B_j$ and fixes $x_0$, so commutes with $a$. These are distinct involutions in the two factors, and their product has order $2$. The generator $b_j$ does not belong to $B_j$, since it changes the base chamber at type $j$. The Coxeter parabolic theorem gives $|a_{x_0}a_{b_jx_0}|=m$; the common exact-label lemma gives $|ab_j|=2m$. Products of old generators have exactly their old orders because $B$ embeds in $G$. Finiteness and the quotient assertion follow from the common theorem and quotient lemma. In the quotient case the faithful automorphism group is obtained by removing a possible chamber kernel, as explained there; no exact Coxeter labels are asserted for an arbitrary quotient.
\end{proof}

In the quotient assertion, parabolics mean intersections of the images of the original maximal parabolics, exactly as in source Definition~2.2. Write $C_F=A_FN/N$ and $C_E=\bigcap_{F\in E}C_F$. Then the old-type parabolic is $C_{X\setminus B_ix_0}\rtimes B_i$, and the new-type parabolic is $C_X\rtimes B$. If the quotient action has a nontrivial chamber kernel $H=C_X$, this latter subgroup is $H\rtimes B$. Only after dividing by $H$ can these intersections be replaced by the subgroups generated by complementary markers. The common quotient lemma proves this reduction; the faithful Coxeter formulas (\ref{eq:onepars}) are unchanged.

The complete group presentation can be written without one generator per element of $X$. Take a presentation of $B$, add $a^2=1$, impose $[a,b_i]=1$ for $i\ne j$, and add
\begin{equation}\label{eq:pairrels}
(a\tau a\tau^{-1})^{\lambda([x_0,\tau x_0])}=1
\end{equation}
for one representative of each finite-labelled diagonal class. The common orbital-presentation proposition proves the converse, so no relation hidden in a selected finite realization is being assumed.

\begin{corollary}[Finite trees]\label{cor:trees}
Every finite tree with Coxeter labels $4$, except possibly one edge with a finite label $k\geq3$, is the diagram of a finite regular hypertope.
\end{corollary}
\begin{proof}
A one-node tree is realized by the rank-one geometry with group $C_2$. A two-node tree of label $k$ is realized by the regular $k$-gon and its dihedral group of order $2k$; use $k=4$ if all labels are $4$. For a larger tree with a distinguished exceptional edge, choose a leaf whose incident edge is not exceptional. Such a leaf exists: the two endpoints of an edge cannot both be leaves in a connected tree with more than two vertices. If no edge is exceptional, choose any leaf. Delete that leaf, apply induction, and let $j$ be its former neighbour. Choose every diagonal label to be $2$. Then $W(D)=C_2^X$, and Theorem~\ref{thm:one} attaches the required leaf with exact label $4$. Its group has finite order
\[
2^{[B:B_j]}|B|.
\]
The subgroup index is at least two by thinness. This constructs the given tree, rather than merely some tree of the same rank.
\end{proof}

\section{Fixed vertices and a second Coxeter graph}
Let $Q$ be a finite Coxeter graph with a distinguished vertex $v$. Write $Y=V(Q)\setminus\{v\}$. Form $K$ on $Y\sqcup X$ as follows: retain $D$ on $X$, retain $Q$ on $Y$, and give every pair $\{y,x\}$, $y\in Y,x\in X$, the label of $\{y,v\}$ in $Q$. An omitted edge means label $2$. This is exactly $\operatorname{Ext}(Q,D)$ with $v$ identified with $x_0$. The action of $B$ fixes $Y$ and acts on $X$. It preserves every label. There is one nontrivial type orbit, $X$, and the other orbits are singletons. The latter have automatic orbit intersection identities; for $X$ use (\ref{eq:admissibility}).

\begin{theorem}\label{thm:fixed}
Put $A=W(K)$. There is a regular hypertope with group $A\rtimes B$, types $Y\sqcup\{*\}\sqcup I$, generators $c_y=a_y,a=a_{x_0},b_i$, and maximal parabolics
\begin{align*}
G_i&=A(Y\cup B_ix_0)\rtimes B_i&& (i\in I),\\
G_*&=A(Y)\rtimes B,\\
G_y&=A((Y\setminus\{y\})\cup X)\rtimes B&& (y\in Y).
\end{align*}
Its diagram is the disjoint union of $Q$ and the diagram of $\beta$, together with one edge between $v$ and $j$, of exact label $2\lambda([x_0,b_jx_0])$; here $v$ denotes the generator $a$. Its finiteness and invariant-quotient variants are exactly those of Theorem~\ref{thm:one}, with $W(K)$ replacing $W(D)$.
\end{theorem}
\begin{proof}
Apply the common theorem to the orbits just described. Its formulas give every displayed subgroup. The generators $c_y$ commute with $B$ because their vertices are fixed. The subgroup generated by $c_y$ and $a$ is the standard Coxeter parabolic on $Y\cup\{x_0\}$, hence has precisely the labels of $Q$. The remaining product orders are those calculated in Theorem~\ref{thm:one}. This proves all edges and all omitted edges, including exactness. The common quotient lemma applies without changing the indexed orbit condition.
\end{proof}

The full presentation adds the fixed generators and their Coxeter relations to (\ref{eq:pairrels}), with $[c_y,b_i]=1$ and the one relation $(c_ya)^{m_{yv}}=1$ when finite. All other mixed relations follow by conjugation, as proved in the common paper. If $Q$ has one vertex, this reduces to Theorem~\ref{thm:one}. If $D$ has one vertex, the graph operation alone gives $Q$, although that degenerate case is outside our hypothesis $|X|>1$ for the geometric theorem.

For clarity, the source's Example~3.11 contains a small mismatch between its text and picture. Its text names only three vertices for the first graph, but the displayed graph has four, in the chain $u_0,u_1,v,w$ with successive labels $3,4,\infty$. We use that displayed graph. The second graph has vertices $h_1,x_0,h_2$ and successive labels $5,3$. Thus $\operatorname{Ext}(Q,D)$ has six vertices: the retained edge $u_0u_1$ has label $3$, each $u_1x$ has label $4$, and each $wx$ has label $\infty$, for $x\in\{h_1,x_0,h_2\}$; the internal second graph remains unchanged, and every other pair has label $2$. This is exactly the first picture on source page~9. Reversing the operation instead copies labels $5$ and $3$ from the second marked vertex to all four vertices of the first graph, while retaining its chain. This is its second six-vertex picture. These explicit lists recover both pictures and explain why the operation is not symmetric.
\section{Classical polytope identifications}
The comparisons in source Corollaries~3.7 and~3.14 use the standard descriptions of the McMullen--Schulte extensions [2, Sections~8B--8C]. We retain those descriptions, with their names and specified marked semidirect products; we do not claim a fresh proof or a new reading of the whole book. Their presentations and maximal-parabolic descriptions are the inputs for the named identifications below. The geometry and presentation of our side have been proved above.

When $\beta$ is a regular $n$-polytope and $j=0$, the classical extension $2^{\beta,D}$ has marked group $W(D)\rtimes B$, with generators $a_{B_0},b_0,\ldots,b_{n-1}$ and the natural vertex action. Omitting $a_{B_0}$ gives $B$. Omitting $b_i$ leaves $B_i$ together with the conjugates of $a_{B_0}$ under $B_i$, hence exactly $A(B_iB_0)\rtimes B_i$. These are (\ref{eq:onepars}). Therefore the identity on the marked group gives an isomorphism of all typed cosets and their incidence relations, proving source Corollary~3.7. With every diagonal label $2$ this is Danzer's $2^\beta$; with the antipodal class labelled $s$ and all other classes labelled $2$ it is the matching construction $2^{\beta,D(s)}$. These are special choices of the same orbit function, not additional hypotheses on Theorem~\ref{thm:one}.

For source Corollary~3.14, let $\delta$ be the universal regular $m$-polytope whose group is the string Coxeter group on $Q$, and select its last vertex $v$. Let $\beta$ again be a regular polytope and take $j=0$. The classical extension denoted $\delta^{\beta,Q}$ in the source uses $W(\operatorname{Ext}(Q,D))\rtimes B$ and the marked generators
\[
c_0,\ldots,c_{m-2},a_{B_0},b_0,\ldots,b_{n-1}.
\]
Omitting $b_i$, $a_{B_0}$ or $c_y$ gives respectively the three displayed families in Theorem~\ref{thm:fixed}: in the last case the remaining $B$-conjugates of $a_{B_0}$ recover all of $X$. Hence again identity on the marked group identifies every coset, type and incidence. Both diagrams are strings; the retained string C-group correspondence identifies these coset geometries with their regular polytopes. This includes the one-vertex $Q$ case. The named notation and its dependence on $D$ follow the source's convention; no uniqueness assertion depending only on an unmarked diagram is intended.

\section{Source scope, repairs and verification}
The complete principal contract of [1] is recovered: both constructions, their invariant quotient statements, the finite tree family, both exact diagram formulas and both classical comparisons. The action and diagonal lemmas are included. The introductory question about all finite-labelled diagrams remains outside the result; neither this paper nor the source proves it. No statement that arbitrary invariant quotients are hypertopes is made: the quotient geometry hypothesis is essential.

Three details of the written derivation are made explicit. First, in source Theorem~3.5 the subgroup subscript must use $(I\setminus M)\cup(I\setminus N)$, not its printed intersection; the preceding product also repeats $M$ where $N$ is required. Equation (\ref{eq:admissibility}) proves the corrected identity. Second, the power equation in source Corollary~3.8 is supplemented by projection to $B$ and the Coxeter rank-two subgroup, proving exact order rather than only divisibility; the same argument supplies Corollary~3.15. Third, the tree induction specifies a leaf whose removal preserves the exceptional edge. These are proof repairs or elaborations of stated results, not resolutions of a newly identified open conjecture. The source's overloaded orbit indices denote intersections over whole orbits, as explicitly clarified in [3, Definition~3.10].

The accompanying finite program tests small marked instances, every subset intersection and the coset flag-transitivity criterion. It is validation, not a premise for any uniform statement above. The exact scopes of the separate Lean identities and the retained Coxeter and classical-polytope inputs are recorded with the source clauses. All other proofs needed for the application are in this paper and its common paper. In particular [3, Corollary~3.17] is reproved in the finite-rank faithful regular case used here, and its full paper is not thereby counted as recovered.
\begin{thebibliography}{9}
\bibitem{PT} C. A. Piedade and P. Tranchida, \emph{Coset geometries acting on their elements: The $j$-diagonals twisting}, arXiv:2608.23223v1, 2026, twelve-page author edition.
\bibitem{MS} P. McMullen and E. Schulte, \emph{Abstract Regular Polytopes}, Cambridge University Press, 2002, Sections~8B--8C; the exact classical construction descriptions are retained as identified in [1].
\bibitem{W} C. A. Piedade and P. Tranchida, \emph{The geometry of wreath and semi-direct products}, arXiv:2605.04646v2, 2026, Definitions~3.8 and~3.10, Corollary~3.17.
\bibitem{common} Human Compression Project, \emph{Coset intersections under group operations}, complete common proof edition, 21 pages.
\end{thebibliography}
\end{document}
